% Build the layer, Chapter 1, second build: a current holder who does the
% wrong thing, and what it then reports.
% Inserted in \S\ref{sec:enforcement} after the two laws of commitments and
% before the proposal to widen the oracle vocabulary (OP-5), which the
% section's last run motivates.
%
% No code or output in this file is typed by hand. \pdexcerpt, \pdrun,
% \pdrunlines and \pdsessionfile print files under excerpts/, which the
% kernel's tools/excerpts.mjs cuts from its source and its recorded runs (see
% figures/pd-listing.tex). To change a listing, change the kernel and cut the
% excerpts again. The prose refers to functions by name and never by line
% number.
\subsection{Build the layer: four shortcuts to a passing suite}
\label{sec:swk-roads}

\noindent\emph{Express lane: an agent told to make a failing test pass has
four ways to get a passing suite that take less work than fixing the bug.
The fence of \S\ref{sec:swk-build} refuses none of them. The guards built
here refuse the first three, each in the transaction that makes the
decision, and a recorded run shows the fourth being admitted.}

The fence refuses a holder whose grant is out of date. It does not refuse
a holder whose grant is current and who makes a harmful change under it, and
it does not check what that holder then reports. This section has given the
theory of both: Theorem~\ref{thm:decidability} for the change, and the two
laws for the report. Here both are built in \pdkernelname{}, in the same
four parts as \S\ref{sec:swk-build}. As there, none of the code or output
below was typed by hand.

The work is the checkout repair again. The suite fails for a real reason.
The module multiplies a price in cents by 1.07 and does not round, so for
1999 it returns 2138.9300000000003, and the test expects 2139. The correct
fix is one call to \texttt{Math.round}. This time the job is assigned to
Drake, and Drake looks for something that takes less work. Four shortcuts
give a passing suite without fixing the bug:

\begin{description}
\setlength{\itemsep}{2pt}
\item[Delete.] Remove the test file. A suite with no tests passes.
\item[Remove the assertion.] Keep the file and remove the one assertion that
fails.
\item[Report success.] Change nothing, and report that the work is done.
\item[Rewrite.] Change the assertion to expect what the code returns.
\end{description}

\noindent Drake tries them in this order and stops at the first that is not
refused.

\subsubsection{Naive: a fence, and a promise closed on the agent's own report}

The obvious design adds a record of promises to the kernel of
\S\ref{sec:swk-build} and changes nothing else. Before it starts, an agent
promises a result: here, that \texttt{shop/checkout.test.ts} will pass. The
kernel stores the promise with the grants and gives it a deadline. When the
agent has finished, it closes the promise, and the record shows it as done.
Each of these steps goes through the kernel's single entry point and is
written to its event log, so anyone can read afterwards who promised what, and
when it was closed.

Two properties of this record are already correct, and the fix keeps them.
The kernel sets the deadline. The request that makes a promise has no field
for a time, and the service refuses a request that contains anything except
the condition, so an agent cannot give itself a year (Law~1,
Property~\ref{prop:law1}). And the promise is not free text. It names a
condition from a short fixed list, a test file that will pass or a claim
that will be released, so that a program can check it later.

The design is naive in its last step. The promise is closed when the agent
that made it asks: the kernel accepts the agent's own report that the
condition holds.

\subsubsection{Break: Drake deletes the test}

Listing~\ref{alg:swk-roads-run-fence} is Drake's run against that design.

\pdrun[caption={The break. Drake deletes the test under a current grant, and
his promise is closed on his own report.},label={alg:swk-roads-run-fence}]{test-deletion-fencing-only}

Drake runs the suite and it fails. He promises that the test will pass,
and the kernel sets the deadline thirty minutes later. He claims the test
file and is granted token 1. It is the newest token for that file, presented
by the agent it was issued to, so the fence admits the deletion at event 4.
The suite now has no test files and passes. At event 6 Drake closes his
promise, and the record of promises shows it as done.

Every step went through the kernel and is in its log, and that is the
situation Theorem~\ref{thm:schneider} describes. A monitor that reads this log
learns of the deletion from event 4, which was written by the transaction
that deleted the file. It can raise an alarm, and someone can restore the
test. It cannot prevent the deletion, because it runs afterwards. The record
of promises fails in a different way. Nothing was hidden from it. The only
evidence it had for the closure was the report of the agent that asked for
it. \emph{Now you try:} Exercise~\ref{ex:swk-roads-monitor} asks what a
monitor could still do here, and by when.

\subsubsection{Fix: a guard for each decision, where the decision is made}

Two decisions in that run were wrong: admitting the deletion, and closing
the promise. Each gets a guard inside the transaction that makes the
decision, which is the remedy of Theorem~\ref{thm:decidability}.

\paragraph{The edit.} Bob's acceptance rule for the repair says that the
checkout test may be edited but not deleted or weakened, so the kernel's
policy lists the file as protected. Listing~\ref{alg:swk-roads-guard} is the
rule in code. \texttt{check\_remove} refuses to delete a protected file.
\texttt{check\_write} compares the current file with the file as the write
would leave it, and refuses a write that leaves fewer assertions than the
current file has. The count comes from \texttt{count\_assertions}, which is
not shown: it counts calls to \texttt{assert} in the text.

\pdexcerpt[caption={The write-path guard. A protected test is not deleted,
and an edit does not lower its count of
assertions.},label={alg:swk-roads-guard}]{ch01b-guard}

The kernel's decision function calls these after the fence has admitted
the grant and before anything is applied, inside the transaction of
Listing~\ref{alg:swk-build-propose}. So a refusal leaves the file unchanged,
and the log records the refusal with its reason. The rule uses only the text
of the file, so the function reads nothing except its arguments, and it
always terminates. Theorem~\ref{thm:decidability} requires exactly that of a
guard. It is also a crude rule: it counts assertions without knowing what
they assert. The Verify part returns to that.

\paragraph{The promise.}
Listing~\ref{alg:swk-roads-close} is the rule for closing, Law~2
(Property~\ref{prop:law2}) in code. \texttt{close} is passed the kernel's
record of the promise and the policy, with the newest grant on the path the
promise is about. It refuses any agent except the one that made the promise,
and it refuses a promise that is already closed. Then \texttt{holds}
decides. For a promise to release a claim, it checks that the promiser no
longer holds a live grant, and the grant is in the state the transaction has
just read. For a promise that a test passes, it calls \texttt{passing},
described below.

\pdexcerpt[caption={\texttt{close}. Only the promiser may ask, and the
condition the promise named must be true.},label={alg:swk-roads-close}]{ch01b-close}

The branch on \texttt{policy.closure} is the only difference between the
naive design and this one. The run in Listing~\ref{alg:swk-roads-run-fence}
was made with the policy set to accept the promiser's own report.

Law~2 requires the condition to be named when the commitment is made, and
building this rule is what decided that. It would be more convenient to let
the agent state at closing what shows that the work is done, and an earlier
wording of the law allowed it. But then a promise about the checkout test
could be closed with any true fact from the list. Drake could claim and
release \texttt{notes.txt}, a file nobody needs, and present that released
claim. So the condition is fixed when the promise is made, and
\texttt{close} is passed only the promise's number. Anyone who relies on the
promise can see the condition before the work begins, and the agent cannot
replace it afterwards. \emph{Now you try:}
Exercise~\ref{ex:swk-roads-swap} asks for the closure that the earlier wording
would have admitted.

\paragraph{The test.}
A released claim can be checked inside the transaction, because the grants
are rows in the database. A test cannot. It has to be run, and the kernel
allows a test up to 10 seconds. The kernel's one transaction cannot stay
open that long, because every other agent would wait for it. So the kernel
runs the test first and closes the promise afterwards
(Listing~\ref{alg:swk-roads-kernel-close}). That creates the gap of
\S\ref{sec:swk-build} in a new place. There Heron checked its lease and
wrote 31 seconds later. Here the kernel runs a test, and a file can change
before the closure commits.

\pdexcerpt[caption={The kernel's side of closing. The test runs before the
transaction, on a numbered state of the
files.},label={alg:swk-roads-kernel-close}]{ch01b-kernel-close}

The fix is the one that worked for the lease. The kernel numbers the
states of the files. The \emph{revision} is the count of changes to files
that it has admitted so far, computed from the event log. \texttt{\#run}
reads the files and the revision in one transaction, so the number
identifies exactly those files, and it runs the test on that copy in a
separate process. The result is then stored with its number. In
Listing~\ref{alg:swk-roads-passing}, \texttt{passing} refuses a closure for
which no run of the promised test was supplied. It passes the run to
\texttt{fresh}, which compares the revision of the run with the current
revision of the files, inside the closing transaction. Only then does it
check whether the test passed.

\pdexcerpt[caption={\texttt{passing} and \texttt{fresh}. A run of the test
is valid only for the revision it ran on.},label={alg:swk-roads-passing}]{ch01b-passing}

\texttt{fresh} is \texttt{newest} from Listing~\ref{alg:swk-build-newest}
with a different number. There the number showed how old a permission was.
Here it shows how old an observation is. In both, the comparison is made by
the transaction that acts on the result.

\keyidea[A check is valid for the state it ran on]{Some checks cannot run inside
the transaction they guard: they take too long, or they need something
outside the database. Such a check can still be used safely when its result
is stored with the number of the state it ran on, and the transaction compares
that number with the current one before acting. This is what the paragraph
after Theorem~\ref{thm:decidability} calls freshness. It is the fence of
\S\ref{sec:swk-build}, applied to an observation where that one was applied
to a lease.}

\subsubsection{Verify: the same agent, then each guard removed}

Listing~\ref{alg:swk-roads-run-kernel} is the same agent again. The kernel now
has Bob's rule in its policy and checks a promise's condition before it closes
the promise. The agent is the same code.

\pdrun[caption={Drake against the guards. Each shortcut is refused, and the
promise stays open.},label={alg:swk-roads-run-kernel}]{test-deletion-kernel}

The deletion is refused as \texttt{protected-delete}. Drake's next attempt
keeps the file and removes the assertion that fails. The guard counts two
assertions before and one after, and refuses the write as
\texttt{assertions-dropped}. Then he changes nothing and asks to close. The
kernel runs the test at revision 0 and gets the same failure as Drake, so it
refuses. Events 4 to 6 are those refusals, each written by the transaction
that refused. Drake stops. The test is unchanged and the suite still fails.
The record shows the promise as open. If it is still open after
\mbox{\texttt{t=1800000}} it will be listed as overdue. Nothing else
happens: \pdkernelname{} takes no action on an overdue promise.

As in \S\ref{sec:swk-build}, a run that passes is not yet evidence that the
guards are what refused. Table~\ref{tab:swk-roads} lists each shortcut with
the guard that refused it and the result of the same run on a kernel built
without that guard.

\begin{table}[H]
\centering
\caption{The four shortcuts in \pdkernelname{}: the guard that refuses each, and
the same run with that guard removed. [internal]}
\label{tab:swk-roads}
\footnotesize
\renewcommand{\arraystretch}{1.2}
\begin{tabularx}{\textwidth}{@{}>{\raggedright\arraybackslash}p{0.19\textwidth} >{\raggedright\arraybackslash}p{0.2\textwidth} >{\raggedright\arraybackslash}X@{}}
\toprule
\textbf{Shortcut} & \textbf{Refused by} & \textbf{With that guard removed} \\
\midrule
Delete the test & \texttt{check\_remove} & the deletion is admitted and the suite passes with no tests. The promise still does not close: the kernel finds no test file to run \\
Remove the failing assertion & \texttt{check\_write} & the edit is admitted. The test then passes with one assertion, and the promise closes \\
Report success & \texttt{holds} & the promise closes while the test fails \\
Rewrite what the test expects & nothing & Listing~\ref{alg:swk-roads-run-rewrite} \\
\bottomrule
\end{tabularx}
\end{table}

The kernel's tests check the first three rows. They state two properties of
Drake's run: the protected test is unchanged at the end, and no promise is
closed while its test fails. Both must hold with the guards in place, and
the corresponding one must fail with each guard removed.

\texttt{fresh} needs its own run, because nothing in Drake's run changes a
file between a test and a closure. A test of the kernel makes such a change.
Wren fixes the module and asks to close her own promise. The kernel runs the
test, which passes at revision 1. Before the closure commits, a second agent
writes the bug back, and the files are at revision 2. \texttt{fresh}
refuses, because the run was at revision 1. With \texttt{fresh} built to
admit any number, the same sequence closes Wren's promise while the test
fails.

\paragraph{The fourth shortcut.}
Listing~\ref{alg:swk-roads-run-rewrite} is the end of a run in which Drake
does not give up after the third refusal.

\pdrunlines[caption={The fourth shortcut. The test is rewritten to expect what
the bug returns, and every check
passes.},label={alg:swk-roads-run-rewrite}]{test-rewrite-kernel}{25-31,45-56}{25\textendash 31 and 45\textendash 56}

He rewrites the assertion to expect 2138.9300000000003. The file is still
there and still has two assertions, so the write-path guard admits the edit
at event 7. The test now passes, at revision 1, so event 9 admits the
closure. The record shows the promise as done, and the bug is not fixed.

Each guard checked correctly what it checks. One counts assertions and the
other runs the test. Neither knows what the test should expect, and a rule
that reads only the test's text cannot distinguish this edit from an honest
correction of a wrong expectation. This is a specific case of the Key idea
``Receipt checks do not judge quality''. The next subsection asks what a
larger set of conditions could add, and
Exercise~\ref{ex:swk-roads-fourth} asks for a rule that refuses this
shortcut, and for its cost.

\paragraph{A promise at the command line.}
This section's first example of a commitment was a promise to release a
claim within thirty minutes. The transcript below makes that promise with
the \texttt{harbor} command, against the service with its clock fixed at
zero.

\pdsessionfile[a promise to release, refused and then closed]{promise-demo}

Alice's first attempt to close is refused by the grants table: she still
holds token 1. The service wrote the words ``within 1800 s'':
\mbox{\texttt{harbor promise}} has no flag that sets a time. After the
release, the same command closes the promise.

\begin{pdboundary}
\setlength{\parskip}{4pt plus 1pt}
The runs show that each shortcut exists and which guard refuses it. Drake
is a script, so they do not show how often a model uses one. The test the
kernel runs is the agent's own code. It gets an empty environment and 10
seconds, but it runs as the kernel's user with no sandbox. A test file that
exits with status 0 before its assertions run counts as passing. This is
the cooperative rail of \S\ref{sec:swk-monitor} again. A change to any file
increases the revision, so an unrelated write can cause a closure to be
refused, and under continual writes it can be refused repeatedly. Law~2
lists four kinds of condition and \pdkernelname{} builds two: a merged
commit needs a repository, and a policy sub-check needs a policy monitor.
Every promise gets the same deadline. Like the fence, these guards apply
only to requests that go through the kernel.
\end{pdboundary}
